%======================================================================
\chapter{Prerequisites: tensor algebra and continuum mechanics}\label{app:prereq}
%======================================================================
% Added after the review of Oct 2026: "a short prerequisites appendix on
% tensor algebra and continuum mechanics".

\framepart{Outline}
This appendix recalls, in a few pages, the tensor algebra and the small-strain
continuum mechanics used throughout the book. It fixes the notation of the
list on page~\pageref{front:notation} and gives each operation in
Cartesian components. Section~\ref{sec:pre-tensors} treats vectors and
tensors, Section~\ref{sec:pre-calculus} the differential operators, and
Section~\ref{sec:pre-continuum} the kinematics, stresses and equilibrium of a
deformable body. Standard references are~\cite{SalenconPre,GurtinPre,ChadwickPre}.

%======================================================================
\section{Vectors and tensors}\label{sec:pre-tensors}
\index{tensor algebra}%
\paragraph{Components.} In an orthonormal basis $(\ve_1,\ve_2,\ve_3)$ a vector is
$\vu=u_i\ve_i$ and a second-order tensor $\tens{a}=a_{ij}\,\ve_i\otimes\ve_j$, with
summation over repeated indices. A second-order tensor is a linear map of
vectors: $(\tens{a}\vv)_i=a_{ij}v_j$. A fourth-order tensor
$\bbC=C_{ijkl}\,\ve_i\otimes\ve_j\otimes\ve_k\otimes\ve_l$ is a linear map of
second-order tensors: $(\bbC:\eps)_{ij}=C_{ijkl}\varepsilon_{kl}$.

\begin{gbox}{Products and contractions}
\begin{itemize}
  \item scalar product $\vu\cdot\vv=u_iv_i$; \ tensor product
  $(\vu\otimes\vv)_{ij}=u_iv_j$, so that $(\vu\otimes\vv)\vw=(\vv\cdot\vw)\,\vu$;
  \item double contraction $\tens{a}:\tens{b}=a_{ij}b_{ij}=\tr(\tens{a}\T\tens{b})$;
  \item transpose $(\tens{a}\T)_{ij}=a_{ji}$; symmetric part
  $\sym\tens{a}=\tfrac12(\tens{a}+\tens{a}\T)$; skew part $\tfrac12(\tens{a}-\tens{a}\T)$;
  \item trace $\tr\tens{a}=a_{ii}=\tI:\tens{a}$; deviator
  $\dev\tens{a}=\tens{a}-\tfrac13(\tr\tens{a})\tI$, with $\tr\dev\tens{a}=0$.
\end{itemize}
\end{gbox}

\paragraph{Orthogonal decompositions.} A symmetric tensor and a skew tensor
are orthogonal for ``$:$''; so are a spherical tensor $c\tI$ and a deviator.
Hence, for symmetric $\tens{a}$,
\[
  \tens{a}:\tens{a}=\tfrac13(\tr\tens{a})^2+\dev\tens{a}:\dev\tens{a},
\]
the identity behind the split of the elastic energy into a volume and a
shape part (Exercise~\ref{exo:posdef}) and behind the von Mises criterion.

\paragraph{Fourth-order identities.}
\index{identity tensor}%
On symmetric tensors, $\bbI$ with
$I_{ijkl}=\tfrac12(\delta_{ik}\delta_{jl}+\delta_{il}\delta_{jk})$ is the identity,
$\tfrac13\tI\otimes\tI$ the projector on spherical tensors and
$\bbIdev=\bbI-\tfrac13\tI\otimes\tI$ the projector on deviators. They are
orthogonal projectors: $\bbIdev:\bbIdev=\bbIdev$, $\bbIdev:(\tI\otimes\tI)=0$.
An isotropic tensor $\bbC=3K\,\tfrac13\tI\otimes\tI+2\mu\,\bbIdev$ is therefore
inverted by inverting its two eigenvalues: $\bbC^{-1}=\frac{1}{3K}\,\tfrac13\tI\otimes\tI+\frac{1}{2\mu}\,\bbIdev$.

\paragraph{Voigt notation.}
\index{Voigt notation}%
A symmetric second-order tensor has six independent components, and a
fourth-order tensor with the symmetries of elasticity has $21$. They are
stored as a column of six and a symmetric $6\times6$ matrix, with the index
pairs $11,22,33,23,31,12\mapsto1,\dots,6$: $\sigma_I=c_{IJ}\varepsilon_J$ with
the engineering shears $\varepsilon_4=2\varepsilon_{23}$, $\varepsilon_5=2\varepsilon_{31}$,
$\varepsilon_6=2\varepsilon_{12}$. The factor $2$ keeps the energy
$\tfrac12\sig:\eps=\tfrac12\sigma_I\varepsilon_I$. This is the notation of the
tables of moduli of Chapter~\ref{ch:elliptic}, equation~\eqref{eq:voigt}, and of
finite element codes; the Mandel variant, with factors $\sqrt2$, is discussed there.

\paragraph{Eigenvalues and invariants.} A symmetric tensor has three real
eigenvalues and an orthonormal eigenbasis (principal axes). Its principal
invariants are $I_1=\tr\tens{a}$, $I_2=\tfrac12\bigl((\tr\tens{a})^2-\tr\tens{a}^2\bigr)$
and $I_3=\det\tens{a}$; for a deviator one uses $J_2=\tfrac12\dev\tens{a}:\dev\tens{a}$
and $J_3=\det\dev\tens{a}$, the variables of isotropic yield criteria.

%======================================================================
\section{Differential operators}\label{sec:pre-calculus}
\index{gradient}\index{divergence}%
With $(\cdot)_{,j}=\partial(\cdot)/\partial x_j$:
\begin{gbox}{Gradient and divergence in Cartesian components}
\[
  (\nabla u)_i=u_{,i},\qquad (\nabla\vu)_{ij}=u_{i,j},\qquad
  \dive\vu=u_{i,i}=\tr\nabla\vu,\qquad (\dive\tens{a})_i=a_{ij,j},
\]
\[
  \Delta u=\dive\nabla u=u_{,ii},\qquad
  \dive(\tens{a}\T\vv)=(\dive\tens{a})\cdot\vv+\tens{a}:\nabla\vv .
\]
\end{gbox}
The last identity, integrated with the divergence theorem, is the integration
by parts of Theorem~\ref{thm:ibp-tensor}. In cylindrical or spherical
coordinates the components of $\nabla$ and $\dive$ contain extra terms from the
derivatives of the basis vectors; they are tabulated in~\cite{SalenconPre}.

%======================================================================
\section{Small-strain continuum mechanics}\label{sec:pre-continuum}
\paragraph{Kinematics.}
\index{strain}\index{compatibility}%
A body occupies $\Omega$; its points move by the displacement $\vu(\vx)$. When
the gradient is small, $|\nabla\vu|\ll1$, the change of length and angle is
measured by the small-strain tensor
\[
  \eps[\vu]=\tfrac12(\nabla\vu+\nabla\vu\T),\qquad
  \varepsilon_{ij}=\tfrac12(u_{i,j}+u_{j,i}).
\]
$\varepsilon_{11}$ is the relative elongation of a fibre along $\ve_1$,
$2\varepsilon_{12}$ the decrease of the right angle between $\ve_1$ and $\ve_2$,
and $\tr\eps$ the relative change of volume. The skew part of $\nabla\vu$ is an
infinitesimal rotation, and $\eps[\vu]=\vzero$ exactly for the rigid
displacements $\vu=\vect{a}+\vect{b}\times\vx$. Conversely, a symmetric field $\eps$
derives from a displacement in a simply connected domain if and only if it
satisfies the compatibility conditions
$\varepsilon_{ij,kl}+\varepsilon_{kl,ij}-\varepsilon_{ik,jl}-\varepsilon_{jl,ik}=0$.

\paragraph{Stresses.}
\index{stress}\index{Cauchy theorem}%
The force per unit area exerted across a surface element of normal $\vn$ is
the traction $\vt$. Cauchy's theorem states that it is linear in $\vn$:
$\vt=\sig\vn$, $t_i=\sigma_{ij}n_j$, with $\sig$ the Cauchy stress tensor. The
balance of moments makes $\sig$ symmetric. Its normal component
$\vn\cdot\sig\vn$ is the normal stress (positive in tension) and its
tangential part the shear stress.

\begin{gbox}{Equilibrium}
\index{equilibrium!local}%
For a body under body forces $\vf$ (per unit volume), in statics,
\[
  \dive\sig+\vf=\vzero\ \text{in }\Omega,\qquad \sig=\sig\T,\qquad
  \sig\vn=\vt\ \text{on }\partial\Omega .
\]
\end{gbox}
These equations follow from the balance of forces and moments on every
sub-domain, with the divergence theorem. Their weak form is the principle of
virtual power (Section~\ref{sec:el-vw}).

\paragraph{Energy.} In an elastic material the work of the stresses,
$\sig:\dot\eps$ per unit volume, is stored as the free energy
$\psi(\eps)$, with $\sig=\partial\psi/\partial\eps$; for linear elasticity
$\psi=\tfrac12\eps:\bbC:\eps$ (Chapter~\ref{ch:elliptic}). When part of the
work is dissipated, as in plasticity, the second law requires a non-negative
dissipation (Chapter~\ref{ch:pl}).

%======================================================================
\framepart{Summary}
\begin{gbox*}
\begin{itemize}
  \item \textbf{Tensors} are linear maps: $\tens{a}\vv$, $\bbC:\eps$; all
  operations are written in Cartesian components with summation over
  repeated indices.
  \item \textbf{Spherical and deviatoric parts} are orthogonal; isotropic
  fourth-order tensors are combinations of the two projectors
  $\tfrac13\tI\otimes\tI$ and $\bbIdev$, and are inverted eigenvalue by eigenvalue.
  \item \textbf{Small strains} $\eps=\sym\nabla\vu$ vanish exactly on rigid
  displacements.
  \item \textbf{Stresses} are defined by Cauchy's theorem $\vt=\sig\vn$; they
  satisfy $\dive\sig+\vf=\vzero$ and are symmetric.
\end{itemize}
\end{gbox*}

%======================================================================
\framepart{Exercises}

\exercise{Contractions}\label{exo:pre-contr}
Let $\tens{a}=\ve_1\otimes\ve_2$. Compute $\tens{a}\ve_2$, $\tens{a}\ve_1$,
$\tens{a}:\tens{a}$, $\tens{a}:\tens{a}\T$, $\tr\tens{a}$ and $\sym\tens{a}$.
\begin{solution}
$\tens{a}\ve_2=\ve_1$, $\tens{a}\ve_1=\vzero$, $\tens{a}:\tens{a}=1$,
$\tens{a}:\tens{a}\T=0$, $\tr\tens{a}=0$,
$\sym\tens{a}=\tfrac12(\ve_1\otimes\ve_2+\ve_2\otimes\ve_1)$: a pure shear.
\end{solution}

\exercise{Deviator of a uniaxial stress}\label{exo:pre-dev}
Consider the uniaxial stress $\sig=\sigma\,\ve_1\otimes\ve_1$.
Compute its deviator $\dev\sig$ and the scalar $\sqrt{\tfrac32\dev\sig:\dev\sig}$.
\begin{solution}
$\dev\sig=\sigma\,\diag(\tfrac23,-\tfrac13,-\tfrac13)$ and
$\dev\sig:\dev\sig=\tfrac23\sigma^2$, so $\sqrt{\tfrac32\dev\sig:\dev\sig}=|\sigma|$:
the von Mises equivalent stress of uniaxial tension is the tensile stress.
\end{solution}

\exercise{Rigid displacements}\label{exo:pre-rigid}
Show that $\vu=\vect{a}+\vect{b}\times\vx$ has $\eps[\vu]=\vzero$, and compute
$\nabla\vu$.
\begin{solution}
$u_i=a_i+\epsilon_{ijk}b_jx_k$, so $u_{i,l}=\epsilon_{ijl}b_j$, which is skew in
$(i,l)$; its symmetric part vanishes.
\end{solution}

\exercise{Inverting isotropic elasticity}\label{exo:pre-inverse}
Using the projectors, invert $\sig=\lambda(\tr\eps)\tI+2\mu\eps$.
\begin{solution}
$\bbC=(3\lambda+2\mu)\tfrac13\tI\otimes\tI+2\mu\bbIdev$, so
$\eps=\frac{\tr\sig}{3(3\lambda+2\mu)}\tI+\frac{1}{2\mu}\dev\sig$, i.e.\
$\eps=\frac{1}{3K}\,\frac{\tr\sig}{3}\tI+\frac{1}{2\mu}\dev\sig$ with
$K=\lambda+\tfrac23\mu$.
\end{solution}

%======================================================================
\begin{chapterbib}{9}
\bibitem{ChadwickPre} P.~Chadwick, \emph{Continuum Mechanics: Concise Theory and Problems}, Dover, 1999.
\bibitem{GurtinPre} M.~E.~Gurtin, \emph{An Introduction to Continuum Mechanics}, Academic Press, 1981.
\bibitem{SalenconPre} J.~Salen\c{c}on, \emph{Handbook of Continuum Mechanics: General Concepts, Thermoelasticity}, Springer, 2001.
\end{chapterbib}
